\documentclass[12pt,a4paper]{article}

\usepackage[utf8]{inputenc}
\usepackage[T1]{fontenc}
\usepackage[brazil]{babel}
\usepackage[top=3cm,left=3cm,bottom=2cm,right=2cm]{geometry}
\usepackage{mathptmx}
\usepackage{inconsolata}
\usepackage{setspace}
\usepackage{indentfirst}
\usepackage{titlesec}
\usepackage{booktabs}
\usepackage{array}
\usepackage{listings}
\usepackage{xcolor}
\usepackage{xurl}
\urlstyle{same}

\onehalfspacing
\setlength{\parindent}{1.25cm}

\titleformat*{\section}{\normalfont\bfseries\normalsize}
\titleformat*{\subsection}{\normalfont\bfseries\normalsize}

\lstset{
    basicstyle=\ttfamily\footnotesize,
    breaklines=true,
    columns=fullflexible,
    keepspaces=true,
    frame=single,
    rulecolor=\color{black!25},
    backgroundcolor=\color{black!4},
    xleftmargin=0.2cm,
    xrightmargin=0.2cm,
    framexleftmargin=4pt,
    framexrightmargin=4pt,
    framextopmargin=4pt,
    framexbottommargin=4pt,
    literate=
        {á}{{\'a}}1 {é}{{\'e}}1 {í}{{\'i}}1 {ó}{{\'o}}1 {ú}{{\'u}}1
        {Á}{{\'A}}1 {É}{{\'E}}1 {Í}{{\'I}}1 {Ó}{{\'O}}1 {Ú}{{\'U}}1
        {ã}{{\~a}}1 {õ}{{\~o}}1 {Ã}{{\~A}}1 {Õ}{{\~O}}1
        {â}{{\^a}}1 {ê}{{\^e}}1 {ô}{{\^o}}1
        {à}{{\`a}}1 {ç}{{\c{c}}}1 {Ç}{{\c{C}}}1
}

\begin{document}

\section{Questão 07: Sistema de apoio ao diagnóstico baseado em casos}

O protótipo recebe sintomas separados por vírgula, recupera casos semelhantes e pode reutilizar o diagnóstico e o tratamento registrados em um deles. O operador revisa a hipótese antes de gravar a nova experiência. O sistema usa casos fictícios para demonstrar o ciclo de raciocínio baseado em casos (RBC) de Aamodt e Plaza [1].

\subsection{Base de casos e representação}

O arquivo \texttt{casos.json} contém 15 casos iniciais. Cada registro tem identificador, lista de sintomas observados, diagnóstico e tratamento registrado. Na primeira execução, o programa carrega os registros em um banco SQLite. Casos retidos são gravados nesse banco e ficam disponíveis em novas sessões.

\begin{table}[htbp]
    \centering
    \begin{tabular}{>{\raggedright\arraybackslash}p{3.8cm} >{\raggedright\arraybackslash}p{2.0cm} >{\raggedright\arraybackslash}p{8.7cm}}
        \toprule
        \textbf{Diagnóstico} & \textbf{Casos} & \textbf{Exemplos de sintomas representados} \\
        \midrule
        Resfriado comum & C01--C03 & Coriza, espirros, dor de garganta e tosse. \\
        \addlinespace
        Influenza & C04--C06 & Febre, tosse, dor muscular e fadiga. \\
        \addlinespace
        Rinite alérgica & C07--C09 & Espirros, coceira no nariz e olhos lacrimejantes. \\
        \addlinespace
        Enxaqueca & C10--C12 & Dor de cabeça, náusea e sensibilidade à luz. \\
        \addlinespace
        Gastroenterite viral & C13--C15 & Diarreia, vômito e dor abdominal. \\
        \bottomrule
    \end{tabular}
\end{table}

Por exemplo, C04 associa febre, tosse, dor muscular e fadiga ao rótulo Influenza. O tratamento registrado para esse caso fictício inclui repouso, hidratação e avaliação clínica. Os exemplos foram construídos a partir de descrições gerais dos quadros [2].

\subsection{Recuperação por similaridade de Jaccard}

O programa compara conjuntos de sintomas após normalizar espaços, caixa e acentos e remover duplicatas. Dados o conjunto da consulta $Q$ e o conjunto de um caso $C$, a similaridade divide a quantidade de sintomas em comum pela quantidade de sintomas distintos na união.

\[
J(Q, C) = |Q \cap C| / |Q \cup C|
\]

O escore varia de 0, quando os conjuntos são disjuntos, a 1, quando são iguais. O programa ordena os resultados por escore decrescente e pelo número do caso em caso de empate. Exibe até três resultados positivos e inclui os demais casos empatados na primeira posição.

\subsection{Reutilização, revisão e retenção}

\textbf{Reutilização.} Quando os melhores vizinhos concordam no diagnóstico, o programa usa o diagnóstico e o tratamento do primeiro colocado. Empate no topo entre diagnósticos distintos ou ausência de sintomas em comum deixa a decisão para a revisão.

\textbf{Revisão.} O operador pode validar a hipótese, informar um diagnóstico e um tratamento revisados ou descartar a consulta. A validação no protótipo é uma decisão do operador da simulação.

\textbf{Retenção.} Ao confirmar a gravação, o programa associa os sintomas da consulta à solução revisada, cria um identificador e insere o caso no SQLite. As próximas consultas já podem recuperá-lo.

\subsection{Consultas e resultados}

\begin{table}[htbp]
    \centering
    \begin{tabular}{>{\raggedright\arraybackslash}p{4.6cm} >{\raggedright\arraybackslash}p{10.2cm}}
        \toprule
        \textbf{Sintomas informados} & \textbf{Resultado observado} \\
        \midrule
        Febre, tosse e dor muscular. & C04 liderou com 3/4 = 75\%. C05 e C06 ficaram com 2/5 = 40\%. O protótipo sugeriu Influenza e exibiu o tratamento de C04. \\
        \addlinespace
        Coriza e espirros. & C01 e C07 empataram com 2/4 = 50\%, com diagnósticos diferentes. O operador registrou Rinite alérgica na revisão e gravou C16. \\
        \addlinespace
        CORÍZA, espirros, espirros, após reiniciar. & A base foi reaberta com 16 casos. C16 liderou com 2/2 = 100\%. Acentos, caixa e duplicatas não alteraram a comparação. \\
        \addlinespace
        Sintoma fictício ausente da base. & Nenhum vizinho teve escore positivo. O sistema não gerou hipótese automática. \\
        \bottomrule
    \end{tabular}
\end{table}

Na primeira consulta, os três sintomas informados estão em C04. A união contém quatro termos porque C04 também registra fadiga. Na segunda, dois quadros distintos compartilham o mesmo par de sintomas, o que demonstra por que o sistema encaminha o empate para revisão.

\subsection{Análise e execução}

Os testes verificam a recuperação, a revisão e a persistência. A similaridade mede sobreposição de sintomas, sem ponderar gravidade, duração ou exames. O escore de 100\% de C16 indica igualdade entre conjuntos de termos e não estima a probabilidade do diagnóstico. A base reduzida e fictícia limita o protótipo ao uso acadêmico.

Na pasta da questão, execute \texttt{python diagnostico\_rbc.py}. O arquivo Demonstracao.md inclui as entradas e saídas capturadas em duas sessões com um banco separado.

\section*{Referências}

\noindent [1] AAMODT, A.; PLAZA, E. \textbf{Case-Based Reasoning: Foundational Issues, Methodological Variations, and System Approaches}. AI Communications, 1994. Disponível em: \url{https://doi.org/10.3233/AIC-1994-7104}.

\noindent [2] MEDLINEPLUS. \textbf{Resfriado, influenza, rinite alérgica, enxaqueca e gastroenterite}. Disponível em: \url{https://medlineplus.gov/commoncold.html}, \url{https://medlineplus.gov/flu.html}, \url{https://medlineplus.gov/ency/article/000813.htm}, \url{https://medlineplus.gov/migraine.html} e \url{https://medlineplus.gov/gastroenteritis.html}.

\end{document}
